\documentclass[12pt]{article}

\usepackage{sbc-template}
\usepackage{graphicx,url}
\usepackage[utf8]{inputenc}
\usepackage[brazil]{babel}

\sloppy

\title{$titulo}

\author{$autor\inst{1}}

\address{Nome da instituição -- Cidade -- Estado -- Brasil
  \email{seu.email@exemplo.com}
}

\begin{document}

\maketitle

\begin{abstract}
Write here the same summary as below, in English. Most SBC events ask for both versions.
\end{abstract}

\begin{resumo}
Em até 10 linhas: o problema, o que foi feito, como foi feito e o principal resultado. Escreva o resumo por último, quando o artigo estiver pronto.
\end{resumo}

\section{Introdução}

Apresente o tema e por que ele importa. Termine com o objetivo do trabalho em uma frase e com um parágrafo dizendo o que vem em cada seção.

\section{Trabalhos Relacionados}

Mostre o que outros autores já fizeram sobre o tema e o que falta, que o seu trabalho resolve. Por exemplo, os conceitos de interação humano-computador de \cite{barbosa2010ihc}.

\section{Metodologia}

Descreva como o trabalho foi feito: participantes, materiais, etapas e como os dados foram analisados. Outra pessoa deve conseguir repetir o estudo lendo esta seção.

\section{Resultados}

Apresente o que foi encontrado, de preferência com tabelas e figuras, e discuta o que os resultados significam.

\section{Conclusão}

Retome o objetivo, resuma as contribuições, diga as limitações e o que pode ser feito em trabalhos futuros.

\bibliographystyle{sbc}
\bibliography{referencias}

\end{document}
